\documentclass[11pt]{amsart}
\usepackage[T1]{fontenc}
\usepackage{amssymb}
\usepackage[margin=1in]{geometry}
\usepackage[hidelinks]{hyperref}
\newcommand{\code}[1]{\nolinkurl{#1}}
\newcommand{\Hom}{\mathsf{Hom}}
\setlength{\emergencystretch}{2em}
\AddToHook{shipout/foreground}{%
 \put(0.5\paperwidth,-\paperheight+0.55cm){\makebox[0pt]{\tiny\ttfamily
 unit/doc/hcat-segal-homology.tex --- 2026-11-01}}}
\title{Homology and Mayer--Vietoris on the semi-Segal foundation}
\date{1 November 2026}
\begin{document}
\maketitle

We prove homology results conditional on a category $D$ with finite
semi-Segal structure, copowers, and stability. Semi-Segal structure alone
does not give homology. This note separates the assumptions from the
constructions and theorems checked in \code{sag/hcat-segal/All.seg}.

\section{The index type and the selected structure}
The code assumes an indexing type named \code{Cat}. For each index
\code{C : Cat}, it supplies a type \code{objects C}, arrow types
\code{arrows C a b}, and the finite semi-Segal data described below.
It then selects \code{D : Cat} and equips this particular structure
with the additional assumptions for homology. Throughout the note,
$D$ denotes this selected structure, and $D(a,b)$ means
\code{arrows D a b}.

The name \code{Cat} is inherited from the original development.
Here it is only an index type: no types of functors between its indices,
or ambient categorical operations on it, are supplied. It is not
asserted to classify all categories. The indexing arrangement allows
reuse of the existing theorem files; for the homology results one
could instead assume a single finite semi-Segal structure $D$ directly.

The results therefore concern the internal structure of $D$.
They do not construct the infinitely coherent ambient category of
categories considered in the original project, nor derive these
internal semi-Segal inputs from its axioms.

\section{The categorical axioms and their consequences}
We assume a type of objects and, for each pair $a,b$, a type of arrows
$D(a,b)$. The arrow types need not be sets. We supply simplex data through
dimension four, subject to the following filling conditions.
\begin{enumerate}
\item \emph{Triangles.} For composable arrows $f:a\to b$ and $g:b\to c$,
there is a family $A_2(f,g,h)$ of triangles, indexed by arrows $h:a\to c$,
and the type
\[
 \sum_{h:D(a,c)} A_2(f,g,h)
\]
is contractible. We \emph{define} $g\circ f$ to be the third edge of the
chosen centre of this type.
\item \emph{Tetrahedra.} For a compatible tetrahedral boundary with one
specified triangular face missing, the type of completions (the missing
triangle together with a tetrahedron) is contractible. Applied to a
boundary made from the chosen composites, this \emph{constructs} the
associativity path.
\item \emph{Four-dimensional simplices.} For a compatible boundary with
one specified tetrahedral face missing, the type of completions (the
missing tetrahedron together with a four-dimensional simplex) is
contractible. A particular such filling \emph{proves} the pentagon for
the associativity path just constructed.
\end{enumerate}
Contractibility means a chosen completion together with a path from it
to every completion. It retains higher coherence, rather than merely
asserting existence or uniqueness in a set.

Precisely, the implementation uses the inner horns $\Lambda^n_1$ for
$n=2,3,4$. Their missing faces are respectively $02$, $023$, and $0234$.
The boundaries are normalized in stages: triangle data become paths
from selected composites, and tetrahedron data become equalities of
triangular pastings. The required normalization equivalences are proved.
We do not assume or construct an infinite semisimplicial object here.

We also assume identity arrows and two families of triangles on
$(1_a,f,f)$ and $(f,1_b,f)$. The unit laws follow from these triangles.
Thus composition, associativity, the unit laws, and \code{hpent} are
derived; they are not separately supplied equations. All laws are paths
in ordinary dependent type theory, with no judgmental categorical laws.

\section{Additional assumptions for homology}
For one such $D$, we assume the following data and universal properties.
\begin{enumerate}
\item \emph{Copowers by small types.} For every small type $X$ and object
$A$, an object $X\otimes A$ and arrows $\eta_x:A\to X\otimes A$ such that
restriction is an equivalence:
\[
 D(X\otimes A,T)\longrightarrow\prod_{x:X}D(A,T),
 \qquad q\longmapsto(q\circ\eta_x)_{x:X}.
\]
This is a universal property of the supplied maps, not just an arbitrary
equivalence between these types.
\item \emph{A zero object $0$.} Both $D(0,T)$ and $D(T,0)$ are
contractible for every $T$.
\item \emph{Suspensions.} For every $B$, a chosen pushout
$\Sigma B=0\amalg_B0$, including its square and universal property.
\item \emph{Stability.} Every pushout square in $D$ is also a pullback
square. This is stated representably: applying $D(S,-)$ gives a pullback
of types for every $S$.
\item \emph{A distinguished object $G$.} The code calls it
\code{unitD}. No monoidal structure or monoidal-unit property is assumed.
\end{enumerate}
The foundations are ordinary dependent type theory with universes and
the existing \code{Theory} parameters \code{univ}, \code{images}, and
\code{surjections}: univalence and replacement, expressed by small image
factorizations and their surjectivity. These supply the truncations and
sets of components used below. Function extensionality is derived.
No new axioms beyond these foundations and the explicit categorical and
homology inputs are used in the results below.

\section{What is proved: homology and its maps}
For a coefficient object $A$ of $D$, we \emph{define}
\[
 H_n(X;A)=\pi_0\Hom_D(\Sigma^nG,X\otimes A),\qquad n\geq0.
\]
Here $X\otimes A$ is the copower indexed by the type $X$.
The development constructs abelian groups on these sets, including $H_0$.
It proves functoriality of the induced maps. Multiplicativity in positive
degrees is proved in \code{HGroupMaps.seg}, and in degree zero in
\code{StableGroupMaps.seg}. The file \code{HomologyMaps.seg} proves that
homotopic functions induce equal maps and that equivalences of types
induce equivalences on every $H_n$; these are now group isomorphisms in
every nonnegative degree. The equivalence proof applies $\pi_0$ to postcomposition
with the invertible tensor map. Commutative squares of functions induce
commutative squares on homology by functoriality and homotopy invariance.

The new file \code{HomologyUnits.seg} identifies every chosen group unit
with the class of the categorical zero map. In positive degrees this
uses the existing pointed suspension--loop equivalence. In degree zero,
the cone of zero is $(0,0,p\cdot p^{-1})$. Cancelling this loop and then
the two contraction paths shows that the stable loop equivalence sends
zero to the constant loop. Transport of the loop-group unit gives the
claim. Postcomposition preserves the zero class, so the induced maps
preserve units in every degree.

The degree-zero additivity proof uses the existing stable equivalence
\[
 \sigma_{B,E}:D(B,E)\simeq\Omega D(B,\Sigma E).
\]
Write $L_{E,T}:D(\Sigma E,T)\simeq\Omega D(E,T)$ for the suspension
universal-property equivalence. Define the suspension of an arrow
$f:E\to T$ by
\[
 \Sigma f=L_{E,\Sigma T}^{-1}(\sigma_{E,T}(f)).
\]
The theorem \code{StableNaturality.sigma-natural} identifies
$\sigma_{B,T}(f\circ c)$ with the pointed loop map of postcomposition
by $\Sigma f$ applied to $\sigma_{B,E}(c)$.
Its proof normalizes $\sigma(c)$ as precomposition of the suspension
meridian by $c$, and $L(g)$ as postcomposition of the meridian by $g$.
Associativity compares the two orders of composition. Commutativity of
loops in stable mapping types makes the comparison independent of the
chosen pointing paths. On components this is a square of equivalences
with a loop-group homomorphism, proving additivity for the original
degree-zero group structure. No new axioms are used.

\section{What is proved: excision and Mayer--Vietoris}
Excision is a theorem: the copower $(-)\otimes A$ sends a given pushout
square of small types to a pushout square in $D$. The proof uses the
copower universal property and the derived pentagon. Stability then
makes the resulting square a pullback as well, giving the fibre
descriptions used for exactness. Existence of all pushouts of small
types is not an additional premise of this conditional theorem: the
particular pushout and its universal property are supplied.

Let $X$ be the pushout of $B\xleftarrow{i}C\xrightarrow{j}E$, with
maps $u:B\to X$ and $v:E\to X$. Suppress $A$ in the notation.
The unchanged \code{HomologySequence.seg} proves, for every $n\geq1$,
\[
 H_{n+1}B\oplus H_{n+1}E\longrightarrow H_{n+1}X
 \xrightarrow{\partial}H_nC\longrightarrow H_nB\oplus H_nE
 \longrightarrow H_nX.
\]
All arrows are homomorphisms and the three interior positions are exact.
The non-boundary maps are $(i_*,j_*)$ and $u_*-v_*$.
The structural boundary is the same map in both exactness assertions;
adjacent degree blocks have the same difference map.

The new \code{MayerVietorisTail.seg}, together with
\code{MayerVietorisZeroGroups.seg}, joins this sequence to
\[
 H_1B\oplus H_1E\longrightarrow H_1X
 \xrightarrow{\partial_1}H_0C\longrightarrow H_0B\oplus H_0E
 \xrightarrow{u_*-v_*}H_0X.
\]
All arrows are homomorphisms, including the existing connecting map
$\partial_1$, and all three interior positions are exact.
The equalizer statement below is equivalent to exactness at the product
by group cancellation.
The map $\partial_1$ is the direct boundary with source object
$G$. Its incoming difference is definitionally the degree-one
difference above. The existing cone and fibre exactness theorems give
\[
 \operatorname{im}(u_*-v_*)=\partial_1^{-1}(0),\qquad
 \operatorname{im}\partial_1=(i_*,j_*)^{-1}(0,0).
\]
The image of $(i_*,j_*)$ is the equalizer of the two final maps
$u_*\operatorname{pr}_1$ and $v_*\operatorname{pr}_2$.
Images mean merely inhabited fibres, expressed by propositional
truncation; no choice of preimages is asserted.

\section{Why the final boundary is additive}
Work first with a pushout square $B\leftarrow C\to E$ with apex $P$
in $D$, and a test object $S$. The two zero arrows
$B,E\to\Sigma C$, together with the suspension meridian of $C$
and the zero-map identifications, form a cocone. The pushout universal
property constructs an arrow $\beta:P\to\Sigma C$ with that cocone.
This construction is \code{ConnectingArrow.Arrow.connecting}.

Let $\delta:D(\Sigma S,P)\to D(S,C)$ be the \emph{existing} direct
boundary. The comparison theorem proves
\[
 \sigma_{S,C}(\delta(g))
   =\Omega(\beta_*)\bigl(L_{S,P}(g)\bigr),
 \qquad g:D(\Sigma S,P),
\]
where $\beta_*$ is postcomposition with $\beta$, equipped with a
specified path identifying its value at zero with zero.
To prove the comparison, evaluate each pullback cone against the cocone
of $\beta$. The pentagon compares the two orders of composition;
coherence of precomposition with zero identifies the endpoints.
At the constant loop, this evaluation is constant. Consequently its
two endpoint paths agree, so on arbitrary loops the evaluation is
precisely a pointed loop map.

On components, $L$ and $\sigma$ are the equivalences used to define
the original source and target group structures. A pointed loop map
preserves loop concatenation, hence transporting its homomorphism law
along these same equivalences proves
\[
 \partial_1(x+y)=\partial_1(x)+\partial_1(y).
\]
This is \code{BoundaryAdditivity.Types.hom}, also exposed as
\code{MayerVietorisZeroGroups.Types.delta-hom} for the exact tail above.
There are no additional axioms or diagram-compatibility hypotheses.
All coherence here follows from identity types, the derived pentagon,
and the given universal properties.

\section{Boundary naturality: the precise proved statement}
Consider two pushout squares in $D$, vertical arrows
$a_C,a_B,a_E,a_P$, and homotopies
$u_1a_B\simeq a_Pu$ and $v_1a_E\simeq a_Pv$.
For a fixed test object $S$, these homotopies induce a map $N_S$ between
the types of cones on the two cospans after applying $D(S,-)$.
A \emph{coherent map of representable squares}, in the theorem as
implemented, additionally supplies the equality
\[
 \operatorname{cone}_1(a_C\circ c)
   =N_S(\operatorname{cone}(c)),\qquad c:D(S,C).
\]
This is the specified cube compatibility for the diagram map. It is
an input to the naturality theorem, not a new axiom about $D$, and does
not mention a boundary map.

From this data, \code{MayerVietorisNaturality.seg} proves
\[
 \partial_1\,(a_P)_*=(a_C)_*\,\partial
\]
for the direct boundary on components, including the final boundary
into degree zero. The two pointing compatibilities needed for this
argument are proved from the existing zero-map coherence laws.

\code{BoundaryLoopComparison.seg} proves that the structural boundary
used in the positive-degree exact sequence equals the loop of the
direct boundary, transported through the chosen suspension equivalences.
Consequently \code{BoundaryNaturality.seg} proves the corresponding
naturality statement for that structural boundary. Its input is the
same coherent map of representable squares. No comparison with a
separately chosen direct boundary in the next degree is asserted.

\section{What remains open}
The naturality theorems above require the stated cube compatibility;
constructing this compatibility from maps of pushout diagrams of types
is still required for a natural long exact sequence on those diagrams.
There is no terminal surjectivity claim onto $H_0X$.
These are conditional, generalized homology results. We have not
constructed a model of the assumptions or identified this construction
with ordinary homology with specified coefficients. Model construction
is independent work.

The existing additivity and reduced suspension results are retained.
The separate sphere calculations additionally assume suspensions of
small types with their pushout universal properties, and a choice of
desuspension in $D$ with paths $\Sigma(\operatorname{desusp}B)=B$.
Those inputs are not needed for the Mayer--Vietoris theorem above.

\section{A concrete interpretation to construct}
The intended external example is the stable higher category of spectra,
with $G$ the sphere spectrum. The copower is
$X\otimes E=\Sigma^\infty_+X\wedge E$, so the displayed definition of
homology becomes
\[
 H_n(X;E)=\pi_n(\Sigma^\infty_+X\wedge E).
\]
Taking the Eilenberg--MacLane spectrum $EM(M)$ of an abelian group $M$
should identify this with ordinary unreduced homology with coefficients
in $M$. This is the standard external interpretation, not a checked
interpretation of our axiom system; see Lurie, \emph{Higher Algebra},
\S1.4, and Nardin, \emph{Introduction to stable homotopy theory},
Chapter 2, especially Exercise 20 for the ordinary-homology comparison.

There is now a checked first ingredient in
\code{sag/hcat-spectrum/EMSpectrum.seg}. Reusing \code{sag/em}, it
constructs pointed types
\[
 E_0=M,\qquad E_{n+1}=K(M,n+1),\qquad
 \Omega E_{n+1}\simeq_* E_n.
\]
At the first stage the equivalence is the band of $M$-torsors, and
its preservation of addition is checked. At later stages it is the
existing delooping equivalence. These are coefficient
$\Omega$-spectrum data; they do not yet define the higher category of
spectra or its copowers.

The remaining interpretation work is to construct mapping types and
the finite semi-Segal structure for spectra, establish the copower and
stability properties with the required universe sizes, and prove the
comparison with ordinary homology, including its maps. The presheaf,
fibration, universe and univalence constructions in \code{unit/core}
provide semantic infrastructure. They do not by themselves supply these
spectral constructions or a verification of every parameter of this
homology root.

\begin{thebibliography}{2}
\bibitem{HA} J. Lurie, \emph{Higher Algebra}, 2017.
\url{https://www.math.ias.edu/~lurie/papers/HA.pdf}.
\bibitem{Nardin} D. Nardin, \emph{Introduction to stable homotopy theory}.
\url{https://ncatlab.org/nlab/files/Nardin-StableHomotopyTheory.pdf}.
\end{thebibliography}

\section{Validation and preservation}
The homology root checks with 2876 declarations and no failures in a fresh
\code{--opened} run (3.13 seconds). The coefficient-spectrum root checks
with 1001 declarations and no failures (0.14 seconds).
The original 169 \code{hcat} source files are preserved; shared theorem
files are relative links. The continuation uses the horn-derived
pentagon. It does not identify the new groups
with groups formed using the old chosen categorical operations.

\section{Why this may be interesting}
The main point is that a finite amount of categorical coherence suffices
for the homology results proved here. Simplex data through dimension
four supply composition, associativity and their pentagon; together
with the stated universal properties, they support abelian groups,
excision and Mayer--Vietoris in every nonnegative degree. Iterating
suspension increases the homological degree without requiring an
infinite semisimplicial object in the formalization. This establishes
sufficiency of the given finite data, not their minimality.

The additivity of the final boundary makes this more than a construction
of groups and an exact sequence of pointed sets. The boundary defined
by the pullback property is proved compatible with the group structures
defined by suspension. Its comparison with a pointed loop map explains
why these independently defined constructions agree. The coherence
needed for that agreement is proved in ordinary dependent type theory,
using the pentagon derived from a horn filler.

The development also separates two mathematical tasks. A model must
supply the finite semi-Segal data and the copower, zero, suspension and
stability properties; the homology proofs can then be reused. These
model obligations are substantial, but they are explicit. The chosen
object $G$ specifies which maps detect homology, so different choices
of $D$, $G$ and coefficients may give different theories. A verified
spectra model, with $G$ the sphere spectrum and Eilenberg--MacLane
coefficients, would connect this finite foundation to ordinary homology.
That interpretation and the remaining naturality comparison are still
needed. The present contribution is the checked conditional argument
and an explicit account of the coherence it uses.

\end{document}
